\documentclass[conference,compsoc]{IEEEtran}
\usepackage{amsmath,amsfonts,amssymb}
\usepackage{algorithmic}
\usepackage{algorithm}
\usepackage{array}
\usepackage[caption=false,font=footnotesize,labelfont=sf,textfont=sf]{subfig}
\usepackage{textcomp}
\usepackage{stfloats}
\usepackage{url}
\usepackage{verbatim}
\usepackage{graphicx}
\usepackage[nocompress]{cite}
\usepackage{listings}
\usepackage{xcolor}
\usepackage[catalan]{babel}
\usepackage[T1]{fontenc}
\usepackage[utf8]{inputenc}
\usepackage{booktabs}
\usepackage{float}
\usepackage{cuted}

% Catalan names for abstract and index terms
\renewcommand{\abstractname}{Resum}
\renewcommand{\IEEEkeywordsname}{Termes d'índex}

% Abstract and IEEEkeywords: bold label, normal-weight body text
\makeatletter
\def\abstract{\normalfont\@IEEEtweakunitybaselinestretch{1.15}%
    \if@twocolumn
      \@IEEEabskeysecsize\noindent\textbf{\textit{\abstractname}}---\normalfont\@IEEEabskeysecsize\relax
    \else
      \bgroup\par\addvspace{0.5\baselineskip}\centering\vspace{-1.78ex}\@IEEEabskeysecsize\textbf{\abstractname}\par\addvspace{0.5\baselineskip}\egroup\quotation\normalfont\@IEEEabskeysecsize%
    \fi\@IEEEgobbleleadPARNLSP}
\def\IEEEkeywords{\normalfont\@IEEEtweakunitybaselinestretch{1.15}%
    \if@twocolumn
      \@IEEEabskeysecsize\vskip 0.5\baselineskip plus 0.25\baselineskip minus 0.25\baselineskip\noindent
      \textbf{\textit{\IEEEkeywordsname}}---\normalfont\@IEEEabskeysecsize\relax
    \else
      \bgroup\par\addvspace{0.5\baselineskip}\centering\@IEEEabskeysecsize\textbf{\IEEEkeywordsname}\par\addvspace{0.5\baselineskip}\egroup\quotation\normalfont\@IEEEabskeysecsize%
    \fi\@IEEEgobbleleadPARNLSP}
\makeatother

% Configuracio de listings per a Java
\lstset{
    language=Java,
    basicstyle=\ttfamily\footnotesize,
    keywordstyle=\color{blue}\bfseries,
    commentstyle=\color{gray}\itshape,
    stringstyle=\color{red!70!black},
    numbers=left,
    numberstyle=\tiny\color{gray},
    numbersep=5pt,
    frame=single,
    breaklines=true,
    showstringspaces=false,
    tabsize=4,
    captionpos=b
}

\hyphenation{ha-mil-to-nià back-track-ing Warns-dorff re-cur-siu al-go-ris-me al-go-ris-mes con-tro-la-dor mul-ti-pli-ca-ti-va heu-rís-ti-ca}

\begin{document}

\title{Recorregut Hamiltonià amb Backtracking \\ Pràctica 2}

\author{\IEEEauthorblockN{Serafí Nebot Ginard}
\IEEEauthorblockA{Universitat de les Illes Balears (UIB)\\
Grau en Enginyeria Informàtica\\
Algorismes Avançats, Curs 2025--2026}}

\maketitle

\begin{strip}
\begin{abstract}
Aquesta memòria presenta el desenvolupament d'una aplicació Java Swing
que resol el problema del recorregut hamiltonià sobre un tauler d'escacs
compartit per dues peces que alternen torns, amb la restricció que cap
peça no pot estar en posició de capturar l'altra en cap moment del
recorregut. L'algorisme emprat és \textit{backtracking} recursiu amb
la heurística de Warnsdorff per a l'ordenació de moviments. L'aplicació
suporta sis tipus de peces ---quatre clàssiques (cavall, reina, torre,
alfil) i dues personalitzades (mussol i serp)--- amb moviments estàtics
i continus, taulers configurables de $4 \times 4$ fins a $12 \times 12$,
visualització pas a pas amb animació, predicció de temps d'execució i
indicador de dificultat. L'arquitectura segueix el patró
Model-Vista-Controlador (MVC) amb concurrència gestionada mitjançant
\texttt{SwingWorker}.
\end{abstract}

\begin{IEEEkeywords}
recorregut hamiltonià, backtracking, heurística de Warnsdorff,
Model-Vista-Controlador, Java Swing, SwingWorker, peces d'escacs,
concurrència
\end{IEEEkeywords}
\end{strip}

\IEEEpeerreviewmaketitle

%% ================================================================
\section{Introducció}

El problema del recorregut hamiltonià consisteix a trobar un camí que
visiti exactament una vegada cada vèrtex d'un graf~\cite{ref_cormen}.
En el context d'un tauler d'escacs, cada casella és un vèrtex i les
arestes connecten caselles segons les regles de moviment d'una peça
determinada. El cas més conegut és el \textit{Knight's Tour}, on un
cavall ha de visitar totes les caselles d'un tauler $8 \times 8$.

Aquesta pràctica generalitza el problema en dues direccions. Primera,
en lloc d'una sola peça, el tauler és compartit per \textbf{dues peces
que alternen torns}, de manera que cada peça visita aproximadament la
meitat de les caselles. Segona, s'incorpora una \textbf{restricció de
no-captura}: en cap pas del recorregut una peça pot estar en posició
d'atacar l'altra. Aquesta restricció converteix el problema clàssic de
recorregut hamiltonià en un problema de satisfacció de restriccions
(\textit{constraint satisfaction problem}, CSP) amb backtracking.

A més, l'aplicació suporta sis tipus de peces amb models de moviment
fonamentalment diferents ---moviment estàtic (salt fix) i moviment
continu (lliscament fins a obstacle)---, la qual cosa genera una gran
varietat de grafs implícits sobre el mateix tauler i fa que la
dificultat variï dràsticament segons la combinació de peces.

L'objectiu de la pràctica és triple:
\begin{enumerate}
    \item Implementar un algorisme de backtracking recursiu amb
          heurístiques d'optimització per resoldre el problema.
    \item Construir una interfície gràfica amb Java Swing que permeti
          configurar els paràmetres, visualitzar el procés de resolució
          i navegar la solució pas a pas.
    \item Seguir estrictament el patró MVC amb concurrència gestionada
          per \texttt{SwingWorker}, garantint una separació neta entre
          la lògica de negoci i la interfície gràfica.
\end{enumerate}

La memòria s'estructura de la següent manera: la Secció~II descriu els
fonaments teòrics; la Secció~III detalla el disseny i la implementació;
la Secció~IV analitza l'algorisme de backtracking i les heurístiques
emprades; la Secció~V presenta els resultats; i la Secció~VI tanca amb
les conclusions.

%% ================================================================
\section{Fonaments Teòrics}

\subsection{Recorregut Hamiltonià}

Un \textbf{recorregut hamiltonià} (o camí hamiltonià) en un graf
$G = (V, E)$ és una seqüència de vèrtexs $v_1, v_2, \ldots, v_{|V|}$
tal que cada vèrtex apareix exactament una vegada i existeix una aresta
$(v_i, v_{i+1}) \in E$ per a tot $i$~\cite{ref_cormen}. Determinar
l'existència d'un camí hamiltonià és un problema NP-complet en el cas
general, la qual cosa implica que no es coneix cap algorisme que el
resolgui en temps polinòmic per a tots els grafs.

En el cas particular d'un tauler d'escacs $N \times M$, el graf
implícit ve determinat per les regles de moviment de la peça. Per al
cavall sobre un tauler $8 \times 8$, el graf té 64 vèrtexs i
aproximadament 168 arestes, i existeixen solucions per a la majoria de
configuracions. Tanmateix, la generalització a dues peces amb
restriccions de captura redueix dràsticament l'espai de solucions
vàlides.

\subsection{Backtracking Recursiu}

El \textit{backtracking} és una tècnica de cerca exhaustiva que
explora l'espai de solucions de forma sistemàtica, construint solucions
candidates de forma incremental i descartant les branques que violen
les restriccions del problema~\cite{ref_sedgewick}. L'esquema general
és:

\begin{enumerate}
    \item \textbf{Cas base:} si la solució parcial és completa, retorna
          èxit.
    \item \textbf{Generació de candidats:} per a la posició actual, es
          generen tots els moviments vàlids.
    \item \textbf{Validació:} per a cada candidat, es comprova que
          satisfà les restriccions.
    \item \textbf{Recursió:} s'aplica el moviment i es crida
          recursivament.
    \item \textbf{Retrocés:} si la recursió falla, es desfà el moviment
          i es prova el següent candidat.
\end{enumerate}

La complexitat en el pitjor cas és exponencial, ja que pot ser necessari
explorar totes les permutacions possibles. Per a un tauler de
$N \times M$ caselles amb dues peces, el nombre de permutacions potencials
creix de forma factorial amb el nombre de caselles. Això fa imprescindible
l'ús d'heurístiques de poda per reduir l'espai de cerca.

\subsection{Heurística de Warnsdorff}

La heurística de Warnsdorff~\cite{ref_warnsdorff}, proposada originalment
per H.~C.~Warnsdorff el 1823 per al problema del recorregut del cavall,
consisteix a prioritzar els moviments cap a caselles amb \textbf{menor
nombre de sortides vàlides} (\textit{onward moves}). La intuïció és que
les caselles amb poques sortides han de ser visitades aviat, ja que
posteriorment serien inaccessibles i farien fracassar la solució.

Formalment, per a cada casella candidata $c$ es calcula el grau de sortida:
\begin{equation}
\label{eq:warnsdorff}
d(c) = |\{c' \in V : (c, c') \in E \text{ i } c' \text{ no visitada}\}|
\end{equation}
i es trien primer les caselles amb $d(c)$ mínim. Això transforma
l'exploració cega del backtracking en una cerca guiada que, en la
pràctica, redueix dràsticament el nombre de retrocessos necessaris.

En aquesta implementació, la heurística s'aplica mitjançant el mètode
\texttt{getValidMovesSorted()} de la classe \texttt{Board}, que ordena
la llista de moviments vàlids per grau de sortida creixent abans de
passar-la al procés de backtracking.

\subsection{Restricció de No-Captura}

La restricció addicional d'aquesta pràctica és que, en cada pas del
recorregut, cap peça no pot estar en posició de capturar l'altra.
Formalment, si $p_1$ i $p_2$ són les posicions actuals de les dues
peces, amb tipus $t_1$ i $t_2$ respectivament, la restricció exigeix:
\begin{equation}
\label{eq:nocaptura}
\neg\, \text{captura}(p_1, t_1, p_2) \;\wedge\; \neg\, \text{captura}(p_2, t_2, p_1)
\end{equation}

on $\text{captura}(p, t, q)$ indica que una peça de tipus $t$ a la
posició $p$ pot atacar la posició $q$. Aquesta comprovació es fa
\textbf{bidireccionalment} (ambdues peces es comproven com a
atacants) i és la principal font de poda de l'espai de cerca.

\subsection{Moviments Estàtics vs.\ Continus}

Les peces d'escacs presenten dos models de moviment fonamentalment
diferents:

\begin{description}
    \item[Moviment estàtic:] la peça salta a una posició fixa
          determinada per un vector de desplaçament. Cada vector
          genera exactament un moviment candidat. Exemples: cavall,
          mussol, serp.
    \item[Moviment continu:] la peça llisca en una direcció fins a
          trobar la vora del tauler o una casella ocupada. Cada vector
          de direcció genera múltiples moviments candidats (un per cada
          casella lliure al llarg de la direcció). Exemples: reina,
          torre, alfil.
\end{description}

Aquesta distinció afecta tant la generació de moviments com la
detecció de captura, ja que les peces amb moviment continu tenen un
radi d'atac que depèn de la configuració actual del tauler.

\subsection{Patró Model-Vista-Controlador}

El patró MVC~\cite{ref_mvc} separa l'aplicació en tres capes
independents. En aquesta pràctica:

\begin{itemize}
    \item \textbf{Model:} el tauler (\texttt{Board}), el solver
          (\texttt{BacktrackingSolver}), el camí hamiltonià
          (\texttt{HamiltonianPath}), les mètriques
          (\texttt{SolverMetrics}) i els tipus de peça
          (\texttt{PieceType}).
    \item \textbf{Vista:} el marc principal (\texttt{MainFrame}), el
          panell de configuració (\texttt{ConfigPanel}), el tauler
          visual (\texttt{BoardPanel}), els controls
          (\texttt{ControlPanel}) i les estadístiques
          (\texttt{StatsPanel}).
    \item \textbf{Controlador:} l'orquestrador (\texttt{SolverController}),
          la tasca de fons (\texttt{SolverTask}) i el predictor
          de temps (\texttt{TimePredictor}).
\end{itemize}

La concurrència es gestiona amb \texttt{SwingWorker}~\cite{ref_java},
que executa l'algorisme de backtracking en un fil de fons i publica
actualitzacions de progrés al fil EDT (\textit{Event Dispatch Thread})
de forma segura.

%% ================================================================
\section{Disseny i Implementació}

\subsection{Capa del Model}

\subsubsection{PieceType (Enum)}

Defineix els sis tipus de peces com a constants d'un \texttt{enum}.
Cada constant encapsula el nom en català, una abreviatura, el tipus
de moviment (\texttt{STATIC} o \texttt{CONTINUOUS}), els vectors de
moviment i un factor de complexitat per a la predicció de temps.
La Taula~\ref{tab:pieces} resumeix les característiques de cada peça.

\begin{table}[H]
\caption{Tipus de peces, moviments i complexitat.\label{tab:pieces}}
\centering
\begin{tabular}{llccc}
\toprule
Peça & Nom & Tipus & Moviments & Complexitat \\
\midrule
Cavall  & \texttt{KNIGHT} & Estàtic  & 8       & 1.0 \\
Reina   & \texttt{QUEEN}  & Continu  & 8 dir.  & 1.5 \\
Torre   & \texttt{ROOK}   & Continu  & 4 dir.  & 0.7 \\
Alfil   & \texttt{BISHOP} & Continu  & 4 dir.  & 0.8 \\
Mussol  & \texttt{OWL}    & Estàtic  & 4       & 1.2 \\
Serp    & \texttt{SNAKE}  & Estàtic  & 16      & 1.4 \\
\bottomrule
\end{tabular}
\end{table}

El mussol es mou en creu a distància~2, saltant per sobre de caselles
ocupades. La serp combina 8 moviments en L estesa $(\pm 3, \pm 1)$ i
$(\pm 1, \pm 3)$ amb els 8 moviments de rei, donant-li una gran
mobilitat en espais oberts però amb un alt factor de complexitat per
a la cerca. El mètode \texttt{toString()} retorna el nom en català, de
manera que els \texttt{JComboBox<PieceType>} mostren automàticament les
etiquetes correctes sense cap codi addicional.

\subsubsection{Position (Record)}

Registre immutable de Java~16 que representa una coordenada
$(r, c)$ al tauler. Proporciona mètodes per calcular la distància
Manhattan i comprovar adjacència. Com a \textit{record}, genera
automàticament \texttt{equals()}, \texttt{hashCode()} i
\texttt{toString()}, garantint l'ús segur com a clau en mapes i
conjunts.

\subsubsection{Board}

Classe central que representa el tauler $N \times M$. Manté una
matriu \texttt{boolean[][]} de caselles visitades i un comptador
\texttt{visitedCount} per detectar eficientment el cas base (totes
les caselles visitades).

Els mètodes principals són:

\begin{itemize}
    \item \texttt{getValidMoves(Position, PieceType)}: genera la
          llista de moviments vàlids. Per a peces estàtiques, aplica
          cada vector una vegada. Per a peces contínues, estén cada
          direcció fins a trobar la vora o una casella visitada.
    \item \texttt{getValidMovesSorted(Position, PieceType)}: ordena
          els moviments per la heurística de Warnsdorff (grau de
          sortida creixent), tal com s'expressa a
          l'equació~(\ref{eq:warnsdorff}).
    \item \texttt{canCapture(Position, PieceType, Position)}: comprova
          si una peça pot atacar una posició. Per a peces contínues,
          verifica la línia de visió sense obstacles.
\end{itemize}

\subsubsection{HamiltonianPath}

Emmagatzema la seqüència ordenada de posicions visitades i l'índex
de la peça que ha fet cada moviment (0 o 1). Actua com una pila
durant el backtracking (\texttt{addMove}/\texttt{removeLastMove}) i
com a objecte de transferència de dades cap a la vista un cop trobada
la solució.

\subsubsection{SolverMetrics}

Recull estadístiques en temps real: nombre d'iteracions, nombre de
retrocessos, profunditat màxima de recursió i temps transcorregut.
El temps es mesura amb \texttt{System.currentTimeMillis()} i es pot
consultar tant durant l'execució (per a actualitzacions en viu) com
després.

\subsubsection{BacktrackingSolver}

Implementa l'algorisme de backtracking recursiu. El mètode públic
\texttt{solve()} inicialitza l'estat i crida el mètode privat
\texttt{backtrack()}, que segueix l'esquema del Llistat~\ref{lst:backtrack}.

\begin{lstlisting}[caption={Esquema del mètode backtrack.},label=lst:backtrack]
boolean backtrack(curPos, curPiece,
        curIdx, otherPos, otherPiece,
        otherIdx, depth) {
  if (board.visitedCount == totalCells)
    return true; // Cas base
  if (stopped) return false;
  metrics.incrementIterations();
  for (Position next :
       board.getValidMovesSorted(
           curPos, curPiece)) {
    if (canCapture(next, curPiece,
            otherPos, otherPiece))
      continue; // Poda: captura directa
    if (canCapture(otherPos, otherPiece,
            next, curPiece))
      continue; // Poda: captura inversa
    board.setVisited(next);
    solution.addMove(next, curIdx);
    // Alternar torn
    if (backtrack(otherPos, otherPiece,
            otherIdx, next, curPiece,
            curIdx, depth+1))
      return true;
    board.unsetVisited(next);
    solution.removeLastMove();
    metrics.incrementBacktracks();
  }
  return false;
}
\end{lstlisting}

Aspectes clau de la implementació:

\begin{itemize}
    \item \textbf{Alternança de torns:} a cada nivell de recursió,
          els paràmetres ``actual'' i ``altre'' s'intercanvien, de
          manera que les dues peces alternen moviments.
    \item \textbf{Doble comprovació de captura:} es verifica tant
          que la peça que es mou no ataqui l'altra com que l'altra
          no ataqui la nova posició, segons
          l'equació~(\ref{eq:nocaptura}).
    \item \textbf{Cancel·lació cooperativa:} el camp
          \texttt{volatile boolean stopped} es comprova a cada
          iteració, permetent aturar la cerca de forma neta des
          del fil principal.
    \item \textbf{Actualitzacions de progrés:} cada 1\,000
          iteracions, s'invoca un \textit{callback} per actualitzar
          les estadístiques a la interfície.
\end{itemize}

\subsection{Capa del Controlador}

\subsubsection{SolverController}

Orquestra tot el cicle de vida de la resolució. Les responsabilitats
principals són:

\begin{enumerate}
    \item \textbf{Inici:} llegeix la configuració del
          \texttt{ConfigPanel}, valida les posicions inicials, crea
          els objectes del model (\texttt{Board},
          \texttt{BacktrackingSolver}) i llança un
          \texttt{SolverTask} en un fil de fons.
    \item \textbf{Aturada:} estableix el \textit{flag} de
          cancel·lació al solver mitjançant
          \texttt{SolverTask.stopSolver()}.
    \item \textbf{Reinici:} atura qualsevol execució en curs,
          esborra el tauler i restableix l'estat visual.
    \item \textbf{Predicció:} actualitza la predicció de temps i
          el nivell de dificultat cada vegada que la configuració
          canvia.
    \item \textbf{Completat:} gestiona els tres resultats
          possibles: solució trobada, solució inexistent, o execució
          aturada per l'usuari.
\end{enumerate}

\subsubsection{SolverTask (SwingWorker)}

Encapsula l'execució del solver en un \texttt{SwingWorker<Boolean, Void>}.
El mètode \texttt{doInBackground()} executa \texttt{solver.solve()} en
un fil de fons. Les actualitzacions de progrés arriben al fil EDT
mitjançant \texttt{SwingUtilities.invokeLater()}, i el mètode
\texttt{done()} gestiona la finalització amb tractament d'excepcions
(\texttt{CancellationException}, \texttt{InterruptedException}).

\subsubsection{TimePredictor}

Classe utilitària que estima el temps d'execució abans d'iniciar la
cerca, basant-se en un model exponencial empíric:
\begin{equation}
\label{eq:prediccio}
\hat{T}(n) = k \cdot e^{\alpha \cdot n} \cdot \bar{c}_{\text{peça}}
\end{equation}
on $n = \text{files} \times \text{columnes}$, $k = 0{,}001$~s és la
constant base, $\alpha = 0{,}35$ és el factor de creixement, i
$\bar{c}_{\text{peça}}$ és la mitjana dels factors de complexitat de
les dues peces. El resultat es presenta en català amb cinc nivells de
dificultat codificats per color, com mostra la
Taula~\ref{tab:dificultat}.

\begin{table}[H]
\caption{Nivells de dificultat i llindars de temps.\label{tab:dificultat}}
\centering
\begin{tabular}{lrl}
\toprule
Nivell & Llindar & Color \\
\midrule
Molt fàcil   & $<1$ s        & Verd \\
Fàcil        & $1$--$10$ s   & Verd-groc \\
Moderat      & $10$--$60$ s  & Taronja-groc \\
Difícil      & $1$--$5$ min  & Taronja fosc \\
Molt difícil & $>5$ min      & Vermell \\
\bottomrule
\end{tabular}
\end{table}

\subsection{Capa de la Vista}

% TODO: Imatge suggerida - captura de pantalla general de l'aplicació
% amb totes les àrees visibles (config, tauler, controls, estadístiques).
% Nom suggerit: figures/aplicacio-general.png
% Quan tinguis la imatge, descomenta les línies següents:
%
% \begin{figure*}[!h]
%     \centering
%     \includegraphics[width=\textwidth]{figures/aplicacio-general.png}
%     \caption{Vista general de l'aplicació amb les quatre àrees
%     principals: configuració (esquerra), tauler (centre), controls
%     (inferior) i estadístiques (inferior dreta).\label{fig:general}}
% \end{figure*}

\subsubsection{MainFrame}

\texttt{JFrame} principal que compon els quatre sub-panells en un
\texttt{BorderLayout}: \texttt{ConfigPanel} a l'oest dins un
\texttt{JScrollPane}, \texttt{BoardPanel} al centre, i
\texttt{ControlPanel} amb \texttt{StatsPanel} al sud. Gestiona
internament la navegació pas a pas (botons \texttt{<<} i \texttt{>>})
i l'animació amb \texttt{javax.swing.Timer}, i actua com a façana de
delegació perquè el controlador interaccioni amb la vista a través
d'un únic punt d'entrada.

\subsubsection{ConfigPanel}

Panell de configuració amb \texttt{GridBagLayout} que implementa
la interfície \texttt{Scrollable} per permetre desplaçament vertical.
Inclou:
\begin{itemize}
    \item \textit{Spinners} per a les dimensions del tauler ($4$--$12$
          files i columnes).
    \item Selectors de tipus de peça (\texttt{JComboBox<PieceType>})
          amb previsualització d'icona.
    \item \textit{Spinners} per a les posicions inicials, amb
          validació de límits automàtica en canviar les dimensions.
    \item Botons de col·locació interactiva (\texttt{JToggleButton}
          en grup mutuament excloent): l'usuari activa el mode de
          col·locació i fa clic directament sobre una casella del
          tauler per establir la posició inicial d'una peça.
    \item Etiquetes de predicció de temps i nivell de dificultat,
          actualitzades en temps real cada vegada que la configuració
          canvia.
\end{itemize}

% TODO: Imatge suggerida - captura del panell de configuració.
% Nom suggerit: figures/config-panel.png
%
% \begin{figure}[!h]
%     \centering
%     \includegraphics[width=0.7\columnwidth]{figures/config-panel.png}
%     \caption{Panell de configuració amb les seccions de mida del
%     tauler, configuració de peces (tipus, icona, posició, botó de
%     col·locació) i predicció de temps.\label{fig:config}}
% \end{figure}

\subsubsection{BoardPanel}

Panell que renderitza el tauler com una quadrícula de cel·les amb
patró d'escacs (clar/fosc). La quadrícula interna es manté sempre
\textbf{perfectament quadrada}, centrada dins l'espai disponible
mittançant un \textit{layout} nul amb càlcul manual de posició
(la dimensió del costat és $\min(\text{ample}, \text{alt})$).

Cada cel·la dibuixa:
\begin{enumerate}
    \item El fons amb colors que distingeixen cel·les visitades
          (verd clar/fosc) de les no visitades (blanc/marró).
    \item El nombre de seqüència a la cantonada superior esquerra.
    \item La icona de la peça (imatge PNG o text com a alternativa)
          centrada a la cel·la.
    \item Un marc daurat per destacar el pas actual durant la
          navegació.
\end{enumerate}

Les cel·les responen a clics del ratolí per a la col·locació
interactiva de peces. El panell ofereix un mode de previsualització
que mostra les posicions inicials de les dues peces sobre el tauler
buit, amb l'icona real de cada peça carregada des dels recursos del
projecte.

% TODO: Imatge suggerida - comparació tauler buit vs. solució.
% Noms suggerits: figures/tauler-buit.png i figures/tauler-solucio.png
%
% \begin{figure}[!h]
%     \centering
%     \subfloat[Previsualització]{\includegraphics[width=0.45\columnwidth]{figures/tauler-buit.png}}
%     \hfill
%     \subfloat[Solució completa]{\includegraphics[width=0.45\columnwidth]{figures/tauler-solucio.png}}
%     \caption{Tauler amb previsualització de posicions inicials (a)
%     i amb una solució completa amb numeració de passos (b).\label{fig:tauler}}
% \end{figure}

\subsubsection{ControlPanel}

Panell de controls amb dues files: botons d'acció a la fila superior
i un control lliscant de velocitat d'animació a la fila inferior.
Els botons són:
\begin{itemize}
    \item \texttt{Iniciar}, \texttt{Aturar}, \texttt{Reiniciar}:
          gestió del cicle de vida del solver.
    \item \texttt{<<}, \texttt{>>}: navegació pas a pas per la
          solució.
    \item \texttt{Reproduir}/\texttt{Pausar}: animació automàtica
          de la solució.
\end{itemize}

El control lliscant ajusta el retard de l'animació entre 50~ms
(ràpid) i 2\,000~ms (lent), amb actualització en temps real durant
la reproducció. Tots els botons empren icones textuals amb
caràcters Unicode del bloc de formes geomètriques (per exemple,
$\blacktriangleright$ per a reproduir i $\blacksquare$ per a aturar),
evitant dependències d'imatges externes per als controls.

\subsubsection{StatsPanel}

Panell d'estadístiques que mostra en temps real: temps transcorregut
(format \texttt{MM:SS} o \texttt{HH:MM:SS}), nombre d'iteracions,
nombre de retrocessos (\textit{backtracks}), profunditat màxima i
estat del solver amb codificació de color. Inclou un \texttt{Timer}
de Swing que actualitza el temps transcorregut cada segon durant
l'execució.

% TODO: Imatge suggerida - captura de pantalla mostrant una solució
% trobada amb el panell d'estadístiques visible.
% Nom suggerit: figures/solucio-stats.png
%
% \begin{figure}[!h]
%     \centering
%     \includegraphics[width=\columnwidth]{figures/solucio-stats.png}
%     \caption{Solució trobada per a un tauler $6 \times 6$ amb dos
%     cavalls, mostrant les estadístiques d'execució.\label{fig:stats}}
% \end{figure}

%% ================================================================
\section{Anàlisi de l'Algorisme}

\subsection{Complexitat Temporal}

L'algorisme de backtracking explora un arbre de cerca on cada node
correspon a una assignació parcial de posicions. En el pitjor cas,
l'arbre té profunditat $N \cdot M$ i cada node pot tenir fins a $b$
branques, on $b$ depèn del tipus de peça i de la mida del tauler.

Per a un cavall en un tauler $N \times M$, el grau màxim és~8. Per
a una reina, el grau pot arribar a $2(N + M) - 2$. La complexitat
en el pitjor cas és:
\begin{equation}
T_{\text{pitjor}} = O(b^{N \cdot M})
\end{equation}

Tanmateix, la heurística de Warnsdorff redueix dràsticament el nombre
de nodes explorats a la pràctica. Per al cavall sobre un tauler
$8 \times 8$, és habitual trobar la solució en menys de 64 passos
(sense retrocessos) quan la heurística és efectiva.

\subsection{Complexitat Espacial}

L'espai addicional consumit per l'algorisme és:
\begin{equation}
S = O(N \cdot M)
\end{equation}
corresponent a la matriu de caselles visitades i la llista de
posicions del camí. La profunditat de la pila de recursió és com
a màxim $N \cdot M$.

\subsection{Impacte de les Restriccions de Captura}

La restricció de no-captura actua com un filtre sobre els moviments
candidats. L'impacte varia segons el tipus de peça:

\begin{itemize}
    \item \textbf{Peces estàtiques} (cavall, mussol, serp):
          la comprovació de captura és $O(1)$ per moviment, ja que
          cal comprovar si la posició de l'altra peça coincideix
          amb un dels desplaçaments fixos.
    \item \textbf{Peces contínues} (reina, torre, alfil):
          la comprovació requereix verificar la línia de visió,
          que és $O(\max(N, M))$ per moviment en el pitjor cas.
\end{itemize}

\subsection{Predicció de Temps}

El model de predicció empra una funció exponencial
(equació~\ref{eq:prediccio}) calibrada empíricament. La justificació
del model exponencial és que el nombre de nodes de l'arbre de cerca
creix exponencialment amb el nombre de caselles. Els factors de
complexitat per peça reflecteixen la diferència observada en temps
d'execució entre combinacions de peces: la reina ($c = 1{,}5$) genera
molts més retrocessos que la torre ($c = 0{,}7$).

%% ================================================================
\section{Resultats i Discussió}

L'aplicació implementada permet resoldre el problema del recorregut
hamiltonià per a una àmplia varietat de configuracions. Les
observacions principals són:

\begin{itemize}
    \item \textbf{Taulers petits} ($4 \times 4$, $5 \times 5$):
          les solucions es troben en menys d'un segon per a la
          majoria de combinacions i posicions inicials.
    \item \textbf{Taulers mitjans} ($6 \times 6$, $7 \times 7$):
          el temps varia significativament segons el tipus de peça.
          Dues torres troben solucions ràpidament; dues reines poden
          requerir minuts.
    \item \textbf{Taulers grans} ($8 \times 8$ i superiors):
          els temps creixen de forma exponencial i algunes
          configuracions poden no trobar solució en temps raonable.
\end{itemize}

La heurística de Warnsdorff és essencial per a la viabilitat
pràctica de l'algorisme. Sense ella, fins i tot taulers de
$5 \times 5$ poden requerir milions de retrocessos.

La funcionalitat de navegació pas a pas i l'animació amb velocitat
ajustable permeten a l'usuari comprendre visualment com les dues
peces cobreixen el tauler i verificar que la solució és correcta.

La col·locació interactiva de peces facilita l'experimentació amb
diferents posicions inicials i permet observar com la posició de
partida influeix dràsticament en el temps de resolució.

Respecte al disseny, l'arquitectura MVC ha facilitat la separació
entre l'algorisme (model) i la visualització (vista), la gestió
de la concurrència al controlador, i la incorporació de
funcionalitats voluntàries sense modificar el model.

%% ================================================================
\section{Conclusions}

Aquesta pràctica ha permès implementar i analitzar un algorisme de
backtracking recursiu per al problema del recorregut hamiltonià amb
dues peces, aplicant la heurística de Warnsdorff i restriccions de
no-captura.

Des del punt de vista algorísmic, s'han verificat experimentalment
els aspectes fonamentals del backtracking: el creixement exponencial
del temps d'execució amb la mida del tauler, l'efecte dramàtic de
les heurístiques d'ordenació de moviments, i la importància de la
poda per restriccions.

Des del punt de vista del disseny, l'aplicació demostra una
implementació MVC on:

\begin{enumerate}
    \item El model és completament independent de la interfície
          gràfica i es podria reutilitzar en un context diferent.
    \item La concurrència es gestiona amb \texttt{SwingWorker},
          garantint que la interfície roman responsiva durant
          cerques que poden durar minuts.
    \item La cancel·lació cooperativa mitjançant un camp
          \texttt{volatile} permet aturar la cerca de forma neta.
    \item Les funcionalitats de visualització (navegació pas a pas,
          animació, previsualització de peces, col·locació
          interactiva) enriqueixen l'experiència d'usuari sense
          afectar la lògica del model.
\end{enumerate}

Com a treball futur, es podrien explorar heurístiques addicionals
(detecció de components desconnectats, cerca bidireccional),
implementar la cerca en paral·lel, o calibrar el model de predicció
amb dades experimentals per millorar la seva precisió.

%% ================================================================

\begin{thebibliography}{9}
\bibliographystyle{IEEEtran}

\bibitem{ref_cormen}
T. H. Cormen, C. E. Leiserson, R. L. Rivest, and C. Stein,
\textit{Introduction to Algorithms}, 4th~ed. Cambridge, MA, USA:
MIT Press, 2022.

\bibitem{ref_sedgewick}
R. Sedgewick and K. Wayne, \textit{Algorithms}, 4th~ed. Upper Saddle
River, NJ, USA: Addison-Wesley, 2011.

\bibitem{ref_mvc}
E. Gamma, R. Helm, R. Johnson, and J. Vlissides, \textit{Design Patterns:
Elements of Reusable Object-Oriented Software}. Reading, MA, USA:
Addison-Wesley, 1994.

\bibitem{ref_warnsdorff}
H. C. Warnsdorff, \textit{Des Rösselsprunges einfachste und allgemeinste
Lösung}. Schmalkalden, Germany, 1823.

\bibitem{ref_knuth}
D. E. Knuth, \textit{The Art of Computer Programming, Vol. 1:
Fundamental Algorithms}, 3rd~ed. Reading, MA, USA: Addison-Wesley, 1997.

\bibitem{ref_java}
Oracle Corporation, ``The Java Tutorials --- Concurrency,''
[Online]. Available:
\url{https://docs.oracle.com/javase/tutorial/essential/concurrency/}

\bibitem{ref_swing}
J. Elliott, R. Eckstein, M. Loy, D. Wood, and B. Cole, \textit{Java
Swing}, 2nd~ed. Sebastopol, CA, USA: O'Reilly Media, 2002.

\end{thebibliography}

\end{document}
